\documentclass{report}
\usepackage{amsmath,amssymb}
\setlength{\parindent}{0mm}
\setlength{\parskip}{1em}
\begin{document}
\begin{center}
\rule{6in}{1pt} \
{\large Michael Saunders \\
{\bf Experiments with iterative computation of search directions within interior methods for constrained optimization}}

Dept of Management Science & Engineering \\ Huang Engineering Center \\ 475 Via Ortega \\ Stanford University \\ Stanford \\ CA 94305-4121
\\
{\tt saunders@stanford.edu}\\
Santiago Akle\end{center}

Our primal-dual interior-point optimizer PDCO has found many applications
for optimization problems of the form
\[
\min\, \varphi(x) \text{ \ st\ \ } Ax=b, \quad l \le x \le u,
\]
in which $\varphi(x)$ is convex and $A$ is a sparse matrix or a linear
operator. We focus on the latter case and the need for iterative
methods to compute dual search directions from linear systems of the
form
\[
A D^2 \!A^T \Delta y = r, \qquad \text{$D$ diagonal and positive definite}.
\]
Although the systems are positive definite, they do not need to be
solved accurately and there is reason to use MINRES rather than
CG (see PhD thesis of David Fong (2011)). When the original problem
is regularized, the systems can be converted to least-squares
problems and there is similar reason to use LSMR rather than LSQR.
Also, $D$ becomes increasingly ill-conditioned as the interior method
proceeds and there is need for some kind of preconditioning, such as
the partial Cholesky approach of Bellavia, Gondzio and Morini (2011).

We present numerical results on matters such as these.


\end{document}
